\section{Numerické riešenie}
Keďže náš manipulátor sa skladá s 6 článkov, tak v nasledovných rovniciach budeme uvažovať 
výpočet pre týchto 6 článkov. Všeobecne však môžeme uvažovať ľubovoľný počet článkov,
a teda aj ľubovoľný počet premenných. Pseudo programy na výpočet týchto rovníc budú
preto všeobecne vyjadrené.

Pre výpočet budeme potrebovať základné vektory a matice. Budeme uvažovať polohový vektor:
\begin{equation}
\mathbf{p} = \begin{pmatrix} p^x \\ p^y \\ p^z \end{pmatrix}
\end{equation}
a vektor kĺbových premenných:
\begin{equation}
\boldsymbol{\theta} = \begin{pmatrix} \theta_1 \\ \theta_2 \\ \theta_3 \\ \theta_4 \\ \theta_5 \\ \theta_6 \end{pmatrix}
\end{equation}

rotačné matice a translačný vektor na výpočet polohy koncového efektora:
\begin{equation}
R_x(\theta) = \begin{pmatrix} 1 & 0 & 0 \\ 0 & \cos(\theta) & -\sin(\theta) \\ 0 & \sin(\theta) & \cos(\theta) \end{pmatrix}
\end{equation}

\begin{equation}
R_y(\theta) = \begin{pmatrix} \cos(\theta) & 0 & \sin(\theta) \\ 0 & 1 & 0 \\ -\sin(\theta) & 0 & \cos(\theta) \end{pmatrix}
\end{equation}

\begin{equation}
R_z(\theta) = \begin{pmatrix} \cos(\theta) & -\sin(\theta) & 0 \\ \sin(\theta) & \cos(\theta) & 0 \\ 0 & 0 & 1 \end{pmatrix}
\end{equation}

\begin{equation}
T = \begin{pmatrix} 0 \\ 0 \\ L \end{pmatrix}
\end{equation}

\subsection{Priama kinematika}
Priama kinematika nám umožňuje vypočítať polohu koncového efektora na základe zadaných kĺbových premenných.
Medzi jednotlivými súradnicovými systémami platia transformácie:
\begin{align}
p_5 &= R_6(\theta_6)(T_6 + p_6) \\
p_4 &= R_5(\theta_5)(T_5 + p_5) \\
p_3 &= R_4(\theta_4)(T_4 + p_4) \\
p_2 &= R_3(\theta_3)(T_3 + p_3) \\
p_1 &= R_2(\theta_2)(T_2 + p_2) \\
p_0 &= R_1(\theta_1)(T_1 + p_1)
\end{align}

Kinematika koncového efektora manipulátora je potom daná ako (transformácia z $p_6$ do $p_0$):
\begin{equation}
p_0(\theta) = R_1(\theta_1)\Bigg(T_1 + R_2(\theta_2)\Big(T_2 + R_3(\theta_3)\big(T_3 + R_4(\theta_4)\left(T_4 + R_5(\theta_5)\left(T_5 + R_6(\theta_6)(T_6)\right)\right)\big)\Big)\Bigg)
\end{equation}


